% Generated draft: requires mathematical and source review.
\documentclass[11pt]{article}
\usepackage[margin=25mm]{geometry}
\usepackage{fontspec}
\setmainfont{Latin Modern Roman}
\usepackage{amsmath,amssymb,mathtools,unicode-math}
\setmathfont{Latin Modern Math}
\usepackage{enumitem,needspace,xurl,hyperref}
\setlength{\parindent}{0pt}
\setlength{\parskip}{5pt}
\setlength{\emergencystretch}{3em}
\newcommand{\N}{\mathbb{N}}
\newcommand{\Z}{\mathbb{Z}}
\newcommand{\Q}{\mathbb{Q}}
\newcommand{\R}{\mathbb{R}}
\newcommand{\C}{\mathbb{C}}
\newcommand{\F}{\mathbb{F}}
\newcommand{\A}{\mathbb{A}}
\newcommand{\PP}{\mathbb{P}}
\newcommand{\E}{\mathbb{E}}
\newcommand{\eps}{\varepsilon}
\newcommand{\dd}{\,\mathrm{d}}
\begin{document}
\section*{UnsolvedMath ID: 6700032}
\textbf{Source problem code:} catalogue entry
\subsection{Definitions and mathematical statement}
The proposed waist--width inequality is
\[
\operatorname{width}_{n-1}(X)\leq C_n\operatorname{waist}_{n-k+1}(X)
\]
for all complete Riemannian $n$-manifolds $X$, where $C_n$ depends only on dimension as printed. The source does not define these width/waist functionals, the range or quantifier of $k$, or their normalization. Urysohn widths, sweepout waists, volumes and length-normalized variants cannot be silently identified; the exact displayed indices are retained.

The appended discussion defines a contractibility radius. With $B_x(r)$ denoting a metric ball, write
\[
\operatorname{contr}(X,r)=\inf\{R\geq r:\ \forall x\in X,\ B_x(r)\hookrightarrow B_x(R)\text{ is null-homotopic}\},
\]
with value $+\infty$ if the class is empty. Thus contraction is permitted within the larger concentric ball, not required intrinsically in the smaller ball.

The supplied text then makes three additional claims. (A) Complete manifolds with cocompact isometry group have $\operatorname{contr}(X,r)=\infty$ for every $r>0$. (B) If $\operatorname{contr}(X,r_i)\leq r_{i+1}$ for $i=1,\ldots,n$ and $r_1\leq\cdots\leq r_{n+1}<\infty$, then
\[
\operatorname{fil.rad}[X]\geq\frac{r_1}{(n+1)!}.
\]
(C) Universal covers $\widetilde X$ of compact aspherical $n$-manifolds have infinite filling radius and $\operatorname{waist}_m(\widetilde X)=\infty$ for $m=1,\ldots,n$.

Claim (A) is incompatible with the given definition: Euclidean space has cocompact isometry group and its convex balls contract within themselves, so its contractibility radius is finite. No replacement sign or inequality is guessed. The filling-radius/fundamental-class and waist conventions for noncompact spaces in (B)--(C), and their reference [41], are absent. These statements are therefore preserved as source claims requiring correction or contextual verification, not used as proved premises for the main conjecture.
\subsection{Short English statement}
Compare the source's waist and width quantities on every complete manifold. Their definitions and an index range are missing, and the appended contractibility-radius discussion contains an explicit inconsistency that must be resolved.
\subsection{Sources}
\begin{itemize}
\item Local UnsolvedMath registry, permanent ID 6700032, exact source code catalogue entry, ``Waist-Width Inequality''; proposed by Misha Gromov, 2017. Registry status partially solved, classification PARTIAL-PROGRESS.
\item Misha Gromov, \emph{101 Questions, Problems and Conjectures around Scalar Curvature} (2017), printed tag [?32], occurrence 33, PDF page 43, \url{https://www.ihes.fr/~gromov/wp-content/uploads/2018/08/101-problemsOct1-2017.pdf}. Supplied normalized pdf-text extraction accessed July 29, 2026; incomplete author-hosted PDF. Its reference [41] is not reproduced. No PDF was accessed here.
\item The supplied summary names Gromov's ``Isoperimetry of waists'' and a 2003 sphere-waist theorem, related filling-radius/width results, Avvakumov--Karasev distribution-waist work and unspecified $\Z/2$ waist results. It reports partial confirmations but no general complete-manifold comparison. Exact theorem hypotheses and full citations are not provided; no equivalence of these waist conventions is assumed.
\item Provenance flags current status and statement-reuse terms NEEDS\_REVIEW. The trailing section-16 heading in the extract is document context, not an extra mathematical hypothesis.
\end{itemize}
\end{document}
